% Sección de experimentación (se incluye desde informe.tex).
% Las tablas de tablas/*.tex y las figuras de figuras/ se regeneran con:
%   python -m benchmarks.run_all && python -m benchmarks.report_tables

\subsection{Metodología}

Los cuatro experimentos están automatizados en \code{benchmarks/} (\code{python -m benchmarks.run\_all})
y guardan sus resultados en \code{benchmarks/results/} (CSV + PNG). Cada experimento crea archivos nuevos
en una carpeta de trabajo, usa el dataset sintético \code{empleados} generado con semilla fija
(ids únicos 1..$N$ en orden aleatorio, registros de 62 bytes) y mide con contadores independientes
(\code{DiskCounter}) y \code{time.perf\_counter()}, con el recolector de basura desactivado durante las
mediciones. Entorno de la corrida: Python 3.12.13 en macOS-26.3.1-arm64-arm-64bit (10 núcleos), corrida del 2026-09-22 20:28:58
.

Las lecturas y escrituras contadas son \textbf{deterministas} (no dependen de la máquina) y son la métrica
principal. Los tiempos son secundarios: los archivos caben en la caché de páginas del sistema operativo,
de modo que una ``lectura'' cuesta microsegundos y el tiempo refleja sobre todo el costo de CPU de
Python. En un disco real el costo estaría dominado por el número de transferencias.

\paragraph{Reproducibilidad.} Para verificar lo anterior repetimos los experimentos 2, 3 y 4 completos en
una segunda máquina (Windows 11 sobre un AMD64 de 8 núcleos, Python 3.12.1). Todas las columnas de I/O
salen idénticas, incluida la desviación estándar: el full scan lee 1\,539 páginas ($\sigma = 0$), la
búsqueda binaria 11.034 ($\sigma = 0.181$) y el B$^+$ y el hash 3 ($\sigma = 0$); en el Experimento~4
coinciden además el fan-out, la altura, el número de hojas, la I/O de carga y los MB transferidos para
los cuatro tamaños de bloque. Lo único que cambia es el tiempo, de 3 a 10 veces mayor en Windows: la
búsqueda puntual pasa de 0.014--22.7\,ms a 0.051--87.6\,ms según la estructura, un rango del 25\,\% con
el B$^+$ de 109 a 570\,ms, y la carga de los 100\,000 registros del Experimento~4 de 3.7--12.1\,s a
19.9--40.9\,s. Las conclusiones no se mueven: en ambas máquinas el B$^+$ le gana al full scan en tiempo
hasta el 5\,\% de selectividad y lo pierde al 10\,\%.

En el Experimento~1 se mide la \emph{inserción estructural} en todas las organizaciones (sin validar la
unicidad de la PK, porque las claves generadas son únicas por construcción): así se compara lo mismo en
todas, ya que el Heap sin índice no valida y los índices B$^+$/hash se miden por separado. En el motor SQL
la unicidad sí se valida (secciones~\ref{sec:archivos} y~\ref{sec:indices}).

%------------------------------------------------------------------------------------------
\subsection{Experimento 1: costo de inserción masiva}

Se insertan lotes $N\in\{10^3, 10^4, 5\cdot10^4, 10^5, 2.5\cdot10^5, 5\cdot10^5\}$ con $B=4096$. Los
índices B$^+$ y hash son índices únicos sobre \code{id} de un Heap: se reporta el costo total (heap +
índice) y el del índice solo.

\begin{table}[H]
\centering\scriptsize
\begin{tabular}{@{}lrrrrrr@{}}
\toprule
Estructura & $N=1\,000$ & $N=10\,000$ & $N=50\,000$ & $N=100\,000$ & $N=250\,000$ & $N=500\,000$ \\
\midrule
\multicolumn{7}{@{}l}{\emph{Tiempo total (s)}} \\
Heap & 0.01 & 0.10 & 0.43 & 0.94 & 2.46 & 4.72 \\
Sequential sin reorg & 0.01 & 0.12 & 0.63 & 1.31 & 3.22 & 6.52 \\
Sequential con reorg & 0.03 & 0.32 & 1.89 & 4.50 & 12.41 & 24.79 \\
B+ & 0.06 & 0.57 & 4.13 & 7.70 & 22.25 & 42.12 \\
Hash & 0.03 & 0.23 & 1.32 & 2.53 & 5.98 & 12.30 \\
\midrule
\multicolumn{7}{@{}l}{\emph{Escrituras de páginas}} \\
Heap & 1\,001 & 10\,001 & 50\,001 & 100\,001 & 250\,001 & 500\,001 \\
Sequential sin reorg & 1\,016 & 10\,154 & 50\,770 & 101\,539 & 253\,847 & 507\,693 \\
Sequential con reorg & 1\,108 & 10\,918 & 55\,077 & 110\,975 & 273\,965 & 547\,805 \\
B+ & 2\,009 & 20\,090 & 100\,516 & 201\,013 & 502\,889 & 1\,005\,706 \\
Hash & 2\,010 & 20\,069 & 100\,263 & 200\,520 & 502\,059 & 1\,004\,109 \\
\midrule
\multicolumn{7}{@{}l}{\emph{Lecturas de páginas}} \\
Heap & 984 & 9\,846 & 49\,230 & 98\,461 & 246\,153 & 492\,307 \\
Sequential sin reorg & 1\,933 & 19\,933 & 99\,933 & 199\,933 & 499\,933 & 999\,933 \\
Sequential con reorg & 3\,288 & 62\,843 & 429\,335 & 958\,162 & 2\,723\,390 & 5\,946\,007 \\
B+ & 2\,576 & 29\,463 & 148\,987 & 298\,383 & 855\,355 & 1\,851\,665 \\
Hash & 2\,989 & 29\,882 & 149\,364 & 298\,724 & 747\,186 & 1\,494\,366 \\
\midrule
\multicolumn{7}{@{}l}{\emph{I/O por inserción (lecturas + escrituras) / N}} \\
Heap & 1.99 & 1.98 & 1.98 & 1.98 & 1.98 & 1.98 \\
Sequential sin reorg & 2.95 & 3.01 & 3.01 & 3.01 & 3.02 & 3.02 \\
Sequential con reorg & 4.40 & 7.38 & 9.69 & 10.69 & 11.99 & 12.99 \\
B+ & 4.58 & 4.96 & 4.99 & 4.99 & 5.43 & 5.71 \\
Hash & 5.00 & 5.00 & 4.99 & 4.99 & 5.00 & 5.00 \\
\midrule
\multicolumn{7}{@{}l}{\emph{Solo el índice (sin el heap): I/O por inserción}} \\
B+ & 2.60 & 2.97 & 3.01 & 3.01 & 3.45 & 3.73 \\
Hash & 3.01 & 3.01 & 3.01 & 3.01 & 3.01 & 3.01 \\
\bottomrule
\end{tabular}

\caption{Experimento 1: tiempo, escrituras, lecturas e I/O por inserción ($B=4096$).}
\label{tab:exp1}
\end{table}

\begin{figure}[H]
\centering
\begin{subfigure}{0.49\textwidth}\includegraphics[width=\textwidth]{exp1_time.png}\caption{Tiempo total.}\end{subfigure}
\begin{subfigure}{0.49\textwidth}\includegraphics[width=\textwidth]{exp1_writes.png}\caption{Escrituras.}\end{subfigure}
\caption{Experimento 1 en escala log-log: todas las curvas son lineales en $N$ (pendiente 1), es decir,
costo por inserción aproximadamente constante.}
\label{fig:exp1a}
\end{figure}

\begin{figure}[H]
\centering
\includegraphics[width=0.78\textwidth]{exp1_io_per_insert.png}
\caption{I/O por inserción medido frente al costo teórico.}
\label{fig:exp1b}
\end{figure}

\paragraph{Análisis.}
\begin{itemize}
\item \textbf{Heap}: 1.98 I/O por inserción para todo $N$, exactamente el costo teórico de 1R + 1W
  (el 0.02 que falta son las inserciones que abren una página nueva, que cuestan solo 1W: una de cada 65).
  Es la inserción más barata y la más rápida ($\approx 4.8$\,s para 500\,000 registros).
\item \textbf{Sequential sin reorganización}: $\approx 3$ I/O por inserción (1R de la página principal,
  1R + 1W de la cabeza de su overflow). La inserción es barata, pero el archivo degenera: con 500\,000
  registros queda 1 página principal y 7\,692 de overflow, y una búsqueda puntual posterior cuesta en
  promedio 4\,026 lecturas (tabla~\ref{tab:exp1seq}, figura~\ref{fig:exp1c}); es decir, sin reorganizar el
  Sequential se comporta como un Heap con peor inserción.
\item \textbf{Sequential con reorganización}: el costo por inserción crece como $\log_2 M$ (4.4 I/O con
  1\,000 registros, 13.0 con 500\,000): cada inserción paga la búsqueda binaria de su página. Las
  reorganizaciones (35 en total para 500\,000 registros) suman solo $\approx 0.1$ escrituras por
  inserción, porque el archivo crece geométricamente entre una y otra (con fill factor 0.75 cada página
  admite $\approx 35$\,\% más registros antes de desbordar: pasa de 48 a 65). A cambio, al terminar la carga la búsqueda
  cuesta $\lceil\log_2 M\rceil\approx 13$ lecturas.
\item \textbf{B$^+$}: el índice solo cuesta $h$R + 1W: 3.0 I/O por inserción mientras $h=2$ y 3.7 con
  500\,000 claves, cuando la altura pasa a 3 (tabla~\ref{tab:exp1forma}); los splits agregan poco
  (una hoja se divide cada $\approx 290$ inserciones). Es la estructura más lenta en tiempo porque cada
  inserción decodifica y reescribe nodos de hasta 508 claves.
\item \textbf{Hash}: 3.0 I/O por inserción del índice para todo $N$ (2R + 1W), independiente de $N$; la
  profundidad global crece de 2 a 11 y el directorio ocupa 1 a 3 páginas. Los splits y las
  duplicaciones del directorio son poco frecuentes, y \code{crc32} reparte las claves de forma tan
  uniforme que los buckets quedan casi llenos (96\,\% de ocupación con 100\,000 claves).
\end{itemize}

\begin{table}[H]
\centering\small
\begin{tabular}{@{}lrrrrrr@{}}
\toprule
$N$ & \multicolumn{3}{c}{Sequential sin reorg} & \multicolumn{3}{c}{Sequential con reorg} \\
\cmidrule(lr){2-4}\cmidrule(l){5-7}
 & $M$ & overflow & búsqueda & $M$ & overflow (reorgs) & búsqueda \\
\midrule
1\,000 & 1 & 15 & 9.5 & 19 & 0 (8) & 4.4 \\
10\,000 & 1 & 153 & 79.9 & 168 & 7 (17) & 7.4 \\
50\,000 & 1 & 769 & 394.3 & 957 & 0 (25) & 9.9 \\
100\,000 & 1 & 1\,538 & 742.2 & 2\,064 & 0 (28) & 11.0 \\
250\,000 & 1 & 3\,846 & 1\,742.6 & 4\,509 & 8 (32) & 12.2 \\
500\,000 & 1 & 7\,692 & 4\,026.3 & 9\,045 & 22 (35) & 13.1 \\
\bottomrule
\end{tabular}

\caption{Sequential File tras la carga: páginas principales ($M$), páginas de overflow y lecturas
promedio de una búsqueda puntual posterior (100 claves aleatorias).}
\label{tab:exp1seq}
\end{table}

\begin{table}[H]
\centering\footnotesize\setlength{\tabcolsep}{4pt}
\begin{tabular}{@{}lrrrrrr@{}}
\toprule
Estructura & $N=1\,000$ & $N=10\,000$ & $N=50\,000$ & $N=100\,000$ & $N=250\,000$ & $N=500\,000$ \\
\midrule
B$^+$: altura $h$ & 2 & 2 & 2 & 2 & 3 & 3 \\
Hash: profundidad global & 2 & 5 & 7 & 8 & 10 & 11 \\
\bottomrule
\end{tabular}

\caption{Forma final de los índices en el Experimento 1.}
\label{tab:exp1forma}
\end{table}

\begin{figure}[H]
\centering
\includegraphics[width=0.92\textwidth]{exp1_sequential_reorg.png}
\caption{Lo que compra la reorganización: sin ella la búsqueda posterior es lineal en $N$ (la cadena de
overflow crece sin límite); con ella se mantiene en $\lceil\log_2 M\rceil$.}
\label{fig:exp1c}
\end{figure}

%------------------------------------------------------------------------------------------
\subsection{Experimento 2: búsquedas puntuales de igualdad}

Con $N=100\,000$ registros ($B=4096$) se ejecutan 1\,000 búsquedas de claves aleatorias existentes en
cada estructura. Las lecturas incluyen traer el registro por RID en los índices.

\begin{table}[H]
\centering\footnotesize\setlength{\tabcolsep}{4pt}
\begin{tabular}{@{}lrrrrrl@{}}
\toprule
Método & lecturas (media) & $\sigma$ & teórico & ms (media) & $\sigma$ (ms) & parámetros \\
\midrule
Full Scan (Heap) & 1\,539.00 & 0.00 & 1\,539 & 22.504 & 4.136 & P = 1539 páginas \\
Búsqueda binaria (Sequential) & 11.03 & 0.18 & 11.03 & 0.093 & 0.206 & M = 2084 páginas principales \\
B$^+$ (IndexScan) & 3.00 & 0.00 & 3 & 0.039 & 0.016 & h = 2 \\
Hash (IndexScan) & 3.00 & 0.00 & 3 & 0.013 & 0.004 & global depth = 8 \\
\bottomrule
\end{tabular}

\caption{Experimento 2: 1\,000 búsquedas de igualdad sobre 100\,000 registros.}
\label{tab:exp2}
\end{table}

\begin{figure}[H]
\centering
\includegraphics[width=\textwidth]{exp2_equality.png}
\caption{Experimento 2 (escala logarítmica). La raya negra marca el costo teórico.}
\label{fig:exp2}
\end{figure}

\paragraph{Análisis.} Los cuatro métodos coinciden con la teoría:
el \textbf{full scan} lee siempre las $P=1\,539$ páginas ($\sigma=0$), la \textbf{búsqueda binaria} sobre
$M=2\,084$ páginas principales hace en promedio $\log_2 M = 11.03$ sondeos ($\lfloor\log_2 M\rfloor$ o
$\lceil\log_2 M\rceil$ según la clave, de ahí $\sigma=0.18$), el \textbf{B$^+$} de altura $h=2$ cuesta
exactamente $h+1=3$ lecturas y el \textbf{hash} exactamente 3 (directorio + bucket + registro, sin
overflow). El índice reduce el costo en un factor de $\approx 500$ frente al full scan. En tiempo, B$^+$
y hash hacen las mismas 3 lecturas pero el hash es casi 3 veces más rápido: busca la clave en el bucket con
una búsqueda de bytes en C, mientras que el B$^+$ decodifica dos nodos completos con \code{struct}. La
búsqueda binaria es $\approx 250$ veces más rápida que el full scan a pesar de no usar índice.

%------------------------------------------------------------------------------------------
\subsection{Experimento 3: búsquedas por rango con selectividad variable}
\label{sec:exp3}

Con $N=100\,000$ ($B=4096$) se consultan rangos \code{id BETWEEN a AND a+k-1} que devuelven el 0.1\,\%,
1\,\%, 5\,\%, 10\,\% y 25\,\% de la tabla (10 rangos aleatorios por selectividad).

\begin{table}[H]
\centering\scriptsize
\begin{tabular}{@{}lrrrrrrrrrr@{}}
\toprule
 & & \multicolumn{3}{c}{B$^+$} & \multicolumn{3}{c}{Sequential} & \multicolumn{3}{c}{Full Scan (Heap)} \\
\cmidrule(lr){3-5}\cmidrule(lr){6-8}\cmidrule(lr){9-11}
Selectividad & $k$ & lecturas & teórico & ms & lecturas & teórico & ms & lecturas & teórico & ms \\
\midrule
0.1\,\% & 100 & 102.3 & 102.0 & 1.64 & 13.2 & 13.1 & 0.25 & 1\,539.0 & 1\,539.0 & 24.14 \\
1\,\% & 1\,000 & 1\,005.5 & 1\,005.0 & 4.05 & 31.8 & 31.9 & 1.55 & 1\,539.0 & 1\,539.0 & 25.05 \\
5\,\% & 5\,000 & 5\,018.8 & 5\,019.0 & 20.60 & 115.2 & 115.2 & 7.47 & 1\,539.0 & 1\,539.0 & 29.76 \\
10\,\% & 10\,000 & 10\,037.2 & 10\,036.0 & 42.06 & 219.5 & 219.4 & 17.34 & 1\,539.0 & 1\,539.0 & 35.86 \\
25\,\% & 25\,000 & 25\,087.8 & 25\,087.0 & 105.57 & 531.7 & 531.9 & 36.17 & 1\,539.0 & 1\,539.0 & 51.48 \\
\bottomrule
\end{tabular}

\caption{Experimento 3: lecturas medidas (media de 10 rangos), costo teórico y tiempo en ms.
Teórico: B$^+$ $= h + \lceil k/f\rceil - 1 + k$; Sequential $= \log_2 M + k/r$ con $r=48$; full scan $=P$.}
\label{tab:exp3}
\end{table}

\begin{figure}[H]
\centering
\includegraphics[width=\textwidth]{exp3_range.png}
\caption{Experimento 3: lecturas y tiempo frente a la selectividad (log-log). La línea vertical marca
el cruce teórico entre el B$^+$ no agrupado y el full scan ($k\approx P$).}
\label{fig:exp3}
\end{figure}

\paragraph{Análisis.}
\begin{itemize}
\item \textbf{B$^+$ (no agrupado)}: cuesta $\approx k$ lecturas (una por registro, porque cada RID apunta
  a una página distinta del Heap) más $h$ y las hojas recorridas (1 hoja cada $\approx 290$ entradas).
  Supera al full scan cuando $k > P$, es decir, con selectividad mayor a $P/N = 1.5$\,\%: con 25\,\%
  hace 25\,088 lecturas, 16 veces más que el full scan. Este es el cruce que justifica que un
  optimizador basado en costos (modo \code{cost} de nuestro planner) elija el full scan para rangos
  amplios.
\item \textbf{Sequential (agrupado)}: la búsqueda binaria ubica la primera página ($\log_2 M\approx 11$
  lecturas) y luego lee páginas contiguas con 48 registros cada una: $k/48$ páginas. Es el mejor método en
  todas las selectividades (532 lecturas para el 25\,\%, un tercio del full scan) y el costo medido
  coincide con la fórmula a la décima.
\item \textbf{Full scan}: $P=1\,539$ lecturas constantes; solo crece el tiempo de CPU (filtrar y
  decodificar más filas).
\item \textbf{Tiempo vs I/O}: en tiempo el B$^+$ le gana al full scan hasta $\approx 5$\,\% aunque ya lea
  3 veces más páginas, porque sus lecturas aleatorias caen en la caché del sistema operativo y cuestan
  poco; en un disco real, donde una lectura aleatoria cuesta de 0.1\,ms (SSD) a varios ms (HDD), el cruce
  se acerca al que predice el conteo de I/O (1.5\,\%).
\end{itemize}

%------------------------------------------------------------------------------------------
\subsection{Experimento 4: sensibilidad al tamaño de bloque}

Para $B\in\{1024, 2048, 4096, 8192\}$ se insertan 100\,000 registros en un Heap con un B$^+$ único sobre
\code{id} y luego se hacen 1\,000 búsquedas puntuales. Como referencia se reporta también el Sequential
(reorganizado) y el hash con el mismo $B$.

\begin{table}[H]
\centering\small
\begin{tabular}{@{}lrrrr@{}}
\toprule
$B$ (bytes) & 1\,024 & 2\,048 & 4\,096 & 8\,192 \\
\midrule
registros por página del heap & 16 & 32 & 65 & 131 \\
páginas del heap ($P$) & 6\,250 & 3\,125 & 1\,539 & 764 \\
entradas por hoja ($L$) & 100 & 202 & 407 & 816 \\
fan-out del interno ($m+1$) & 125 & 253 & 509 & 1\,021 \\
altura $h$ medida & 3 & 3 & 2 & 2 \\
altura teórica [mín, máx] & [3, 3] & [3, 3] & [2, 3] & [2, 2] \\
hojas / internos & 1\,422 / 17 & 714 / 5 & 341 / 1 & 161 / 1 \\
I/O de inserción (heap + B$^+$) & 591\,248 & 563\,984 & 499\,396 & 499\,046 \\
\quad de los cuales en el B$^+$ & 397\,497 & 367\,108 & 300\,934 & 299\,809 \\
lecturas por búsqueda ($h+1$) & 4.00 & 4.00 & 3.00 & 3.00 \\
I/O total (inserción + 1000 búsquedas) & 595\,248 & 567\,984 & 502\,396 & 502\,046 \\
MB transferidos (I/O $\times B$) & 581.3 & 1\,109.3 & 1\,962.5 & 3\,922.2 \\
tiempo de inserción (s) & 3.70 & 5.44 & 7.53 & 12.08 \\
Sequential: $M$ / lecturas por búsqueda & 8\,334 / 13.04 & 4\,167 / 12.03 & 2\,084 / 11.04 & 1\,021 / 10.00 \\
Hash: profundidad global / buckets & 10 / 1\,024 & 9 / 512 & 8 / 256 & 7 / 128 \\
\bottomrule
\end{tabular}

\caption{Experimento 4: efecto del tamaño de página ($N=100\,000$). La altura teórica es el intervalo
entre nodos llenos y nodos al 50\,\%.}
\label{tab:exp4}
\end{table}

\begin{figure}[H]
\centering
\includegraphics[width=\textwidth]{exp4_block_size.png}
\caption{Experimento 4: fan-out y altura del B$^+$ (izq.) e I/O totales y bytes transferidos (der.).}
\label{fig:exp4}
\end{figure}

\begin{figure}[H]
\centering
\includegraphics[width=0.62\textwidth]{exp4_search_reads.png}
\caption{Lecturas por búsqueda puntual según $B$: B$^+$ ($h+1$), Sequential ($\log_2 M$) y hash (3).}
\label{fig:exp4b}
\end{figure}

\paragraph{Análisis.} El fan-out crece linealmente con $B$ (125, 253, 509 y 1\,021) y la altura baja de
3 a 2 al pasar de 2\,048 a 4\,096 bytes, siempre dentro del intervalo teórico. Como cada búsqueda cuesta
$h+1$, pasa de 4 a 3 lecturas, y como cada inserción cuesta $h$R + 1W en el índice, las I/O de inserción
del B$^+$ bajan de 397\,497 a 299\,809 (las del Heap, 2 por registro, no dependen de $B$). El total de
I/O baja un 16\,\% de $B=1024$ a $B=8192$, pero el volumen transferido (I/O $\times B$) crece casi
linealmente: de 581\,MB a 3\,922\,MB. Duplicar $B$ más allá de 4\,096 no reduce la altura para
100\,000 claves (la raíz ya cabe en una página) y solo encarece cada transferencia: páginas de 4\,KB,
el tamaño de página del sistema operativo, son el mejor compromiso para este volumen. El Sequential
muestra el mismo efecto ($\log_2 M$ baja de 13 a 10 lecturas) y el hash mantiene 3 lecturas con
cualquier $B$, con una profundidad global que baja de 10 a 7 porque los buckets son más grandes.
